\documentclass[11pt,a4paper]{article}
\usepackage[utf8]{inputenc}
\usepackage[T1]{fontenc}
\usepackage{lmodern}
\usepackage[portuguese]{babel}
\usepackage[a4paper,margin=2.5cm]{geometry}
\usepackage{xcolor}
\definecolor{TextEspresso}{HTML}{4A4238}
\definecolor{NudeTaupe}{HTML}{BFAFA6}
\definecolor{SoftBlush}{HTML}{D9C5B2}
\usepackage{titlesec}
\usepackage{enumitem}
\usepackage{listings}
\usepackage{hyperref}
\hypersetup{colorlinks=true, linkcolor=TextEspresso, urlcolor=TextEspresso}

\titleformat{\section}{\color{TextEspresso}\Large\bfseries}{\thesection}{0.5em}{}[{\color{NudeTaupe}\titlerule[0.8pt]}]
\titleformat{\subsection}{\color{TextEspresso}\large\bfseries}{\thesubsection}{0.5em}{}

\lstset{
  basicstyle=\ttfamily\small,
  breaklines=true,
  frame=single,
  rulecolor=\color{NudeTaupe},
  columns=fullflexible,
  showstringspaces=false
}

\newcommand{\guiaotitle}[2]{%
  \begin{center}
  {\Large\bfseries\color{TextEspresso} #1}\\[0.3cm]
  {\normalsize\color{NudeTaupe} #2}\\[0.2cm]
  {\color{NudeTaupe}\rule{4cm}{0.5pt}}
  \end{center}
  \vspace{0.5cm}
}

\begin{document}
\guiaotitle{Guião 04 -- Docker: Primeiros Contentores}{Infraestruturas de Computação na Nuvem, Aula 04}

\section{Objetivos}
\begin{itemize}
  \item Observar diretamente as primitivas de isolamento do kernel (namespaces e cgroups) usadas por um contentor.
  \item Correr contentores a partir de imagens públicas do Docker Hub.
  \item Escrever um Dockerfile simples e construir uma imagem própria.
  \item Usar volumes para persistir dados fora do ciclo de vida do contentor.
  \item Criar uma rede Docker do tipo bridge e ligar dois contentores através dela.
  \item Comparar um contentor de aplicação (Docker) com um contentor de sistema (LXC do Proxmox).
\end{itemize}

\section{Pré-requisitos e Material}
\begin{itemize}
  \item A VM \texttt{icn-lab} no Proxmox, criada no Guião 01, com o Docker já instalado e acessível por SSH.
  \item Acesso à Internet a partir da VM, para descarregar imagens do Docker Hub.
  \item Editor de texto (\texttt{nano}, \texttt{vim}) dentro da VM, para escrever o Dockerfile e o código da aplicação.
  \item Acesso à interface web do Proxmox, para a comparação final com LXC.
\end{itemize}

\section{Introdução}
Nesta aula prática vamos explorar Docker enquanto ferramenta de gestão de contentores. Todo o trabalho decorre \textbf{dentro da VM \texttt{icn-lab}} preparada no Guião 01 -- ou seja, vamos correr contentores dentro de uma máquina virtual, que por sua vez corre sobre o hypervisor KVM do servidor Proxmox. Esta sobreposição de camadas é intencional e comum em infraestruturas reais de cloud.

Começamos por observar concretamente os namespaces e cgroups estudados na teoria, passamos pelo ciclo \texttt{pull} / \texttt{run} / \texttt{stop}, escrevemos o nosso próprio Dockerfile, e terminamos com volumes, redes personalizadas e uma comparação direta com o contentor LXC visto na Aula 03.

\section{Passo a Passo}
\begin{enumerate}[label=\arabic*.]

  \item \textbf{Verificar o Docker na VM.} Estabeleça a ligação SSH à sua VM e confirme que o Docker instalado no Guião 01 está operacional:
\begin{lstlisting}
$ ssh utilizador@IP_da_VM
$ docker --version
$ docker run hello-world
\end{lstlisting}
  O último comando descarrega uma pequena imagem de teste e confirma que o daemon está a funcionar corretamente.

  \item \textbf{Observar os namespaces de um contentor.} Arranque um contentor de longa duração e compare a sua visão do sistema com a da VM que o hospeda:
\begin{lstlisting}
$ docker run -d --name observar alpine sleep 3600
\end{lstlisting}
  Compare a árvore de processos vista de dentro e de fora do contentor:
\begin{lstlisting}
# Dentro do contentor (namespace PID proprio):
$ docker exec observar ps aux

# Fora, na VM (o mesmo processo, outro PID):
$ ps aux | grep sleep
\end{lstlisting}
  Registe o PID que o processo \texttt{sleep} tem \emph{dentro} do contentor e o PID que tem \emph{na VM}. Esta diferença é o namespace PID em ação.

  \item \textbf{Listar os namespaces atribuídos ao contentor.} Obtenha o PID do processo no host e inspecione os seus namespaces:
\begin{lstlisting}
$ PID=$(docker inspect -f '{{.State.Pid}}' observar)
$ sudo ls -l /proc/$PID/ns
\end{lstlisting}
  Identifique na listagem os namespaces \texttt{pid}, \texttt{net}, \texttt{mnt}, \texttt{uts} e \texttt{ipc} estudados na aula teórica. Compare com os do processo \texttt{init} da própria VM:
\begin{lstlisting}
$ sudo ls -l /proc/1/ns
\end{lstlisting}
  Os identificadores numéricos que diferem correspondem aos recursos efetivamente isolados.

  \item \textbf{Observar o isolamento de rede e de hostname.}
\begin{lstlisting}
$ docker exec observar hostname
$ docker exec observar ip addr show
$ hostname
$ ip addr show
\end{lstlisting}
  O contentor tem hostname e interfaces de rede próprios -- namespaces UTS e NET.

  \item \textbf{Aplicar e verificar limites com cgroups.} Corra um contentor com memória limitada e confirme o limite efetivamente aplicado pelo kernel:
\begin{lstlisting}
$ docker run -d --name limitado --memory=256m --cpus=0.5 alpine sleep 3600
$ docker stats --no-stream
\end{lstlisting}
  Confirme o limite diretamente no sistema de ficheiros de cgroups da VM:
\begin{lstlisting}
$ PID2=$(docker inspect -f '{{.State.Pid}}' limitado)
$ cat /proc/$PID2/cgroup
\end{lstlisting}
  Consoante a versão de cgroups em uso, o valor do limite de memória pode ser consultado em \texttt{/sys/fs/cgroup/...} (v2, hierarquia unificada) ou em \texttt{/sys/fs/cgroup/memory/...} (v1). Verifique qual a versão ativa na sua VM:
\begin{lstlisting}
$ stat -fc %T /sys/fs/cgroup
\end{lstlisting}
  O resultado \texttt{cgroup2fs} indica cgroups v2; \texttt{tmpfs} indica v1.

  \item \textbf{Limpar os contentores de observação.}
\begin{lstlisting}
$ docker rm -f observar limitado
\end{lstlisting}

  \item \textbf{Correr um contentor nginx.} Descarregue e execute a imagem oficial do nginx, publicando a porta 80 do contentor na porta 8080 do host:
\begin{lstlisting}
$ docker run -d --name web-teste -p 8080:80 nginx
\end{lstlisting}
  Confirme no navegador (\texttt{http://localhost:8080}) que o servidor responde. Liste os contentores em execução e pare o contentor no final:
\begin{lstlisting}
$ docker ps
$ docker stop web-teste
$ docker rm web-teste
\end{lstlisting}

  \item \textbf{Explorar um contentor alpine interativo.} Alpine é uma imagem base minimalista, muito usada como ponto de partida para outras imagens:
\begin{lstlisting}
$ docker run -it --rm alpine sh
/ # cat /etc/os-release
/ # exit
\end{lstlisting}
  A opção \texttt{--rm} remove automaticamente o contentor assim que este termina.

  \item \textbf{Escrever uma pequena aplicação.} Crie uma pasta \texttt{app-teste/} com um ficheiro \texttt{app.py} com uma aplicação Flask muito simples:
\begin{lstlisting}
from flask import Flask
app = Flask(__name__)

@app.route("/")
def hello():
    return "Ola a partir do meu contentor!\n"

if __name__ == "__main__":
    app.run(host="0.0.0.0", port=5000)
\end{lstlisting}
  E um ficheiro \texttt{requirements.txt}:
\begin{lstlisting}
flask
\end{lstlisting}

  \item \textbf{Escrever o Dockerfile.} No mesmo diretório, crie um ficheiro \texttt{Dockerfile}:
\begin{lstlisting}
FROM python:3.12-alpine
WORKDIR /app
COPY requirements.txt .
RUN pip install --no-cache-dir -r requirements.txt
COPY app.py .
EXPOSE 5000
CMD ["python", "app.py"]
\end{lstlisting}

  \item \textbf{Construir a imagem.} A partir do diretório \texttt{app-teste/}:
\begin{lstlisting}
$ docker build -t minha-app:1.0 .
$ docker images
\end{lstlisting}

  \item \textbf{Correr a imagem localmente.}
\begin{lstlisting}
$ docker run -d --name minha-app -p 5000:5000 minha-app:1.0
$ curl http://localhost:5000
\end{lstlisting}
  Confirme que a resposta corresponde ao texto definido em \texttt{app.py}.

  \item \textbf{Usar um volume para persistência.} Crie um volume nomeado e monte-o num contentor, para que os dados sobrevivam à remoção do contentor:
\begin{lstlisting}
$ docker volume create dados-teste
$ docker run -it --rm -v dados-teste:/dados alpine sh
/ # echo "conteudo persistente" > /dados/ficheiro.txt
/ # exit
$ docker run -it --rm -v dados-teste:/dados alpine cat /dados/ficheiro.txt
\end{lstlisting}
  Note que o ficheiro criado no primeiro contentor está disponível no segundo, mesmo o primeiro já tendo sido removido.

  \item \textbf{Criar uma rede bridge personalizada.}
\begin{lstlisting}
$ docker network create rede-teste
$ docker network ls
\end{lstlisting}

  \item \textbf{Ligar dois contentores através da rede personalizada.} Corra dois contentores nessa rede e confirme que conseguem comunicar pelo nome:
\begin{lstlisting}
$ docker run -d --name servidor --network rede-teste nginx
$ docker run -it --rm --network rede-teste alpine sh
/ # wget -O- http://servidor
/ # exit
\end{lstlisting}
  Ao contrário da rede bridge por omissão, uma rede personalizada permite resolução de nomes entre contentores pelo nome do contentor.

  \item \textbf{Comparar com um contentor de sistema (LXC).} Na aula anterior criámos um contentor LXC no Proxmox. Compare os dois modelos:
\begin{lstlisting}
# No contentor Docker (aplicacao): um unico processo
$ docker run -d --name comparar nginx
$ docker exec comparar ps aux
\end{lstlisting}
  Se ainda tiver (ou recriar) um contentor LXC no Proxmox, observe a diferença:
\begin{lstlisting}
# No node Proxmox, dentro do contentor LXC (sistema):
$ pct enter <ctid>
root@icn-ct:~# ps aux
root@icn-ct:~# systemctl list-units --type=service | head
\end{lstlisting}
  Registe quantos processos correm em cada caso. O contentor Docker tem tipicamente um processo (mais os seus filhos diretos); o LXC tem um sistema de init completo com vários serviços.

  \item \textbf{Limpeza final.}
\begin{lstlisting}
$ docker rm -f servidor minha-app comparar
$ docker network rm rede-teste
$ docker volume rm dados-teste
\end{lstlisting}

\end{enumerate}

\section{Entregável / Questões de Verificação}
\begin{enumerate}[label=\alph*)]
  \setlength\itemsep{0.2cm}
  \item Registe o PID do processo \texttt{sleep} visto \emph{dentro} do contentor e \emph{na VM}. Explique a diferença com base no namespace PID.
  \item Compare a saída de \texttt{ls -l /proc/<pid>/ns} para o contentor e para o processo 1 da VM. Que namespaces diferem, e o que significa cada um deles estar isolado?
  \item Qual a versão de cgroups ativa na sua VM, e como a determinou? Que limite de memória foi efetivamente aplicado ao contentor \texttt{limitado}?
  \item Qual a diferença entre a imagem \texttt{minha-app:1.0} e o contentor criado a partir dela?
  \item Porque é que o ficheiro criado no volume \texttt{dados-teste} sobreviveu à remoção do primeiro contentor?
  \item Explique, por palavras suas, porque é que dois contentores na rede bridge por omissão do Docker não se conseguem descobrir mutuamente pelo nome, ao contrário do que acontece numa rede personalizada.
  \item Quantos processos observou no contentor Docker e quantos no contentor LXC? Relacione com a distinção entre contentor de aplicação e contentor de sistema.
  \item Nesta aula, um contentor Docker corre dentro de uma VM KVM, que corre sobre o servidor Proxmox. Enumere as camadas de virtualização envolvidas e indique, para cada uma, se partilha ou não o kernel com a camada inferior.
  \item Envie o \texttt{Dockerfile} utilizado, juntamente com uma captura de ecrã do comando \texttt{curl} sobre \texttt{localhost:5000} com a resposta da aplicação.
\end{enumerate}

\end{document}
